% Wersja polska (robocza) main_v4.tex, 23.11.2026; wiążąca jest angielska.
\documentclass[10pt,twocolumn]{article}
\usepackage[margin=0.75in]{geometry}
\usepackage[T1]{fontenc}
\usepackage[utf8]{inputenc}
\usepackage[polish]{babel}
\usepackage{amsmath,amssymb,amsthm}
\usepackage{booktabs}
\usepackage{multirow}
\usepackage{xcolor}
\usepackage{tikz}
\usepackage[hidelinks]{hyperref}

\newtheorem{proposition}{Stwierdzenie}
\newtheorem{lemma}{Lemat}
\newtheorem{definition}{Definicja}
\newcommand{\todo}[1]{\textcolor{red}{[TODO: #1]}}
\newcommand{\N}{\mathcal{N}}
\newcommand{\Pa}{\mathcal{P}}
\newcommand{\ncr}{\operatorname{cr}}

\title{Czy zamiany powinny się krzyżować? Sąsiedztwa zamian złożonych od Motzkina do inwolucji dla permutacyjnego problemu przepływowego}
\author{Zgłoszenie anonimowe}
\date{}

\begin{document}
\maketitle

\begin{abstract}
Ruch złożony w permutacyjnym problemie przepływowym stosuje naraz kilka
zamian końców przedziałów i jest wybierany dokładnym algorytmem, który
minimalizuje sumę zmian długości uszeregowania wywołanych przez
poszczególne zamiany. Dynasearch dopuszcza tylko zamiany na przedziałach
rozłącznych, a sąsiedztwo Motzkina dopuszcza także zamiany zagnieżdżone.
Uzupełniamy tę oś. Rodzina $\N_c$ dopuszcza co najwyżej $c$ par
krzyżujących się zamian, w $\Pa_1$ każda zamiana może krzyżować się
z co najwyżej jedną inną, a największym członkiem rodziny jest zbiór
wszystkich inwolucji. Liczymy te sąsiedztwa za pomocą ułamka łańcuchowego
i pokazujemy, że ograniczona liczba krzyżowań mnoży liczbę Motzkina
jedynie przez czynnik wielomianowy, $|\N_c|\sim M_n (n/3)^c/c!$.
Najlepszy ruch znajdujemy dokładnie: dla $c=0$ programowaniem dynamicznym
po przedziałach, dla $c=1$ i $\Pa_1$ programowaniem dynamicznym po
rdzeniach z dwóch łuków, a dla inwolucji skojarzeniem o maksymalnej wadze.
Wszystkie rodziny korzystają z tych samych zamian-kandydatów i ich delt,
więc wybór ruchu dokłada niewiele do kosztu $O(mn^3)$ obliczania delt.
Ograniczeniem jest addytywność. Dwie poprawiające zamiany na przedziałach
rozłącznych realizują większą część przewidywanego wspólnego zysku niż
dwie zamiany zagnieżdżone lub krzyżujące się. Przy wyborze addytywnym
jakość iterowanego przeszukiwania lokalnego powinna więc maleć wzdłuż tej
osi, a najlepszym wyborem powinny być zamiany rozłączne. Na 120
instancjach Taillarda iterowane przeszukiwanie lokalne z budżetem
$10$\,s potwierdza taką kolejność w każdej klasie, a Dynasearch wyprzedza
najlepszą pojedynczą zamianę tylko nieznacznie. Dwie kontrole wskazują
przyczynę. Gdy każdy ruch zachowuje tylko dwie najbardziej poprawiające
zamiany, wszystkie rodziny działają podobnie, więc krzyżowania szkodzą
przez oddziaływanie wielu jednoczesnych zamian. Gdy ruch jest wybierany
przez QUBO, które uwzględnia zmierzone oddziaływania par, realizuje on
niemal cały przewidywany zysk, a ruchy z krzyżowaniami dorównują
Dynasearchowi.
\end{abstract}

% ---------------------------------------------------------------------
\section{Wstęp}

Permutacyjny problem przepływowy z kryterium długości uszeregowania,
$F|prmu|C_{\max}$, jest NP-trudny dla $m\geq 3$
maszyn~\cite{garey1976} i stanowi standardowe pole testowe dla
przeszukiwania lokalnego~\cite{taillard1993,ruiz2007}. Dynasearch~\cite{congram2002}
pokazał, że wykładniczy zbiór niezależnych zamian parami można przeszukać
dokładnie programowaniem dynamicznym, co czyni ruch złożony jednostką
przeszukiwania w ramach przeszukiwania sąsiedztw o bardzo dużym
rozmiarze~\cite{ahuja2002}. W problemie przepływowym stosuje się dwa
sąsiedztwa złożone z zamian końców przedziałów. Dynasearch łączy zamiany,
których przedziały są parami rozłączne. Sąsiedztwo Motzkina~\cite{motzkin_own}
dopuszcza także zamiany zagnieżdżone i jest liczone liczbami Motzkina.

Oba sąsiedztwa zabraniają jednej konfiguracji: dwóch zamian $(a,b)$
i $(c,d)$ z $a<c<b<d$, które krzyżują się po narysowaniu ich jako łuków
nad permutacją. Sam ruch tego nie wymaga. Każdy zbiór zamian o parami
różnych końcach jest inwolucją i wyznacza poprawnego sąsiada, a najlepsza
inwolucja przy celu addytywnym jest skojarzeniem o maksymalnej wadze,
które można znaleźć w czasie wielomianowym. Zakaz krzyżowań istnieje
dlatego, że program dynamiczny po przedziałach nie potrafi ich wyrazić.
Ten artykuł pyta, czy ten zakaz coś kosztuje: \emph{czy zamiany w ruchu
złożonym powinny móc się krzyżować?}

Odpowiadamy rodziną sąsiedztw uporządkowaną według tego, jak bardzo dwie
zamiany mogą się nakładać. Zaczyna się od Dynasearcha, przechodzi przez
sąsiedztwo Motzkina i przez ruchy z ograniczoną liczbą krzyżowań, a kończy
na wszystkich inwolucjach. Cała rodzina używa jednego zbioru zamian-kandydatów
i jednego zbioru delt, więc jej członkowie różnią się wyłącznie tym, jakie
podzbiory kandydatów wolno im wybierać. Pozwala to oddzielić wpływ
nakładania się od wpływu kosztu. Metaheurystykę, jej parametry
i protokół eksperymentalny z wcześniejszych prac~\cite{motzkin_own,windowed_own}
zostawiamy bez zmian, więc każdą nową liczbę można odczytać na tle
opublikowanych.

Artykuł wnosi cztery wyniki. Po pierwsze, definiuje rodziny z krzyżowaniami
$\N_c$ i $\Pa_r$ oraz je liczy. Ułamek łańcuchowy daje dokładne rozmiary,
a lemat typu Poissona pokazuje, że stała liczba krzyżowań zostawia tempo
wzrostu na poziomie $3^n$, podczas gdy wszystkie inwolucje rosną szybciej
niż każda funkcja wykładnicza. Po drugie, podaje dokładne algorytmy wyboru
dla każdego poziomu osi i pokazuje, że są tanie w porównaniu z deltami,
z których korzystają. Po trzecie, mierzy, jak addytywność zależy od
geometrii pary zamian: rozłączne, zagnieżdżone lub krzyżujące się. Ten
pomiar wyjaśnia, dlaczego większe sąsiedztwa tego rodzaju więcej
przewidują, a mniej dostarczają. Po czwarte, ocenia rodziny na
benchmarku Taillarda przy regule wyboru z wcześniejszych prac
i podaje, co ruchy złożone dodają do najlepszej pojedynczej zamiany.
Osobny wariant, który jednokrotnie sprawdza rzeczywistą długość
uszeregowania wybranego ruchu, pokazuje, jak dużą część straty to
sprawdzenie usuwa. Spojrzenie na rodzinę przez QUBO dodaje piąty punkt.
Największe sąsiedztwo na osi ma najrzadsze QUBO, a tylko QUBO może
przenieść zmierzone oddziaływania par, których żaden program po
przedziałach nie wyrazi. Rozwiązany dokładnie na małym zbiorze kandydatów,
ten model usuwa stratę ruchów z krzyżowaniami.

% ---------------------------------------------------------------------
\section{Tło}

\paragraph{Ruchy złożone w szeregowaniu.} Dynasearch wprowadzono dla
problemu jednej maszyny z sumaryczną ważoną spóźnialnością~\cite{congram2002};
późniejsze prace przyspieszyły jego przeszukiwanie~\cite{ergun2006}
i dodały wstawienia~\cite{grosso2004}. Bo\.zejko i Wodecki analizowali
multiruchy zamian dla problemu przepływowego~\cite{bozejko2007}.
Sąsiedztwo Fibonacciego łączy rozłączne zamiany sąsiednie~\cite{fibonacci_own,bozejko2026sr},
a sąsiedztwo Motzkina łączy zamiany niekrzyżujące się~\cite{motzkin_own}.
W problemie z sumaryczną ważoną spóźnialnością zmiany kosztu niezależnych
zamian się sumują, a program dynamiczny jest dokładny. W problemie
przepływowym tak nie jest, więc każde z tych sąsiedztw optymalizuje
funkcję zastępczą.

\paragraph{Sąsiedztwa przeszukiwane przez skojarzenia.} Sarvanov
i Doroshko~\cite{sarvanov1981} oraz Gutin~\cite{gutin1999} przeszukują
sąsiedztwa o rozmiarze silniowym dla problemu komiwojażera za pomocą
zadania przydziału, a Deineko i Woeginger~\cite{deineko2000} przeglądają
sąsiedztwa wykładnicze dla problemu komiwojażera i kwadratowego problemu
przydziału. Sąsiedztwo inwolucji należy do tej tradycji. Nie zbadano
natomiast osi pomiędzy programem po przedziałach a skojarzeniem ani tego,
czy ruch wzdłuż niej pomaga.

\paragraph{Struktura drugorzędowa RNA.} Ta sama oś istnieje w składaniu
RNA, gdzie pary zasad są łukami nad sekwencją. Algorytm
Nussinova~\cite{nussinov1978} jest programem po przedziałach bez
krzyżowań. Pseudowęzły są krzyżowaniami: Rivas i Eddy~\cite{rivas1999}
oraz Akutsu~\cite{akutsu2000} podają programy dynamiczne dla klas
ograniczonych, ogólne pseudowęzły są NP-trudne w modelach
energetycznych~\cite{lyngso2000}, a Tabaska i in.~\cite{tabaska1998}
otrzymują dowolne pseudowęzły ze skojarzenia o maksymalnej wadze, gdy
ocena jest addytywna. Nasz poziom $c=0$ odpowiada algorytmowi Nussinova,
$c=1$ prostemu pseudowęzłowi, a inwolucje podejściu przez skojarzenia.

\paragraph{Krzyżowania w skojarzeniach.} Rozkład liczby krzyżowań
w skojarzeniach i inwolucjach jest klasyczny~\cite{touchard1952,riordan1975}
i ma postać ułamka łańcuchowego~\cite{flajolet1980}. Krzyżowania
i zagnieżdżenia mają ten sam rozkład~\cite{kasraoui2006}, a skojarzenia
$k$-niekrzyżujące się policzyli Chen i in.~\cite{chen2007}. Wykorzystujemy
te wyniki do oszacowania rozmiarów sąsiedztw.

% ---------------------------------------------------------------------
\section{Rodzina z krzyżowaniami}\label{sec:family}

$n$ zadań jest wykonywanych na maszynach $1,\dots,m$ w kolejności
$\pi=(\pi(1),\dots,\pi(n))$, a $p_{r,j}$ jest czasem wykonania zadania
$j$ na maszynie $r$. Czasy zakończenia spełniają
$C_{r,t}=\max(C_{r-1,t},C_{r,t-1})+p_{r,\pi(t)}$ oraz
$C_{\max}(\pi)=C_{m,n}$.

\begin{definition}
\emph{Ruch elementarny} $s(i,j)$, $i<j$, zamienia miejscami $\pi(i)$
i $\pi(j)$ i zostawia zadania pomiędzy nimi na miejscu. \emph{Ruch
złożony} jest zbiorem $\sigma$ ruchów elementarnych o parami różnych
końcach, czyli inwolucją pozycji; sąsiadem jest $\pi\circ\sigma$. Dwie
zamiany $(a,b)$ i $(c,d)$ z $a<c$ są \emph{rozłączne}, jeśli $b<c$,
\emph{zagnieżdżone}, jeśli $d<b$, i \emph{krzyżujące się}, jeśli $c<b<d$.
$\ncr(\sigma)$ jest liczbą par krzyżujących się w $\sigma$.
\end{definition}

\begin{figure}[t]
\centering
\begin{tikzpicture}[x=0.36cm,y=0.36cm,every node/.style={font=\scriptsize}]
\foreach \o/\lab in {0/rozłączne,7/zagnieżdżone,14/krzyżujące się}{
  \draw[gray] (\o,0)--(\o+5,0);
  \foreach \x in {0,...,5} \fill (\o+\x,0) circle (1.2pt);
  \node at (\o+2.5,-0.9) {\lab};
}
\draw (0,0) to[out=70,in=110] (2,0);
\draw (3,0) to[out=70,in=110] (5,0);
\draw (7,0) to[out=70,in=110] (12,0);
\draw (8.5,0) to[out=70,in=110] (10.5,0);
\draw (14,0) to[out=70,in=110] (17,0);
\draw (16,0) to[out=70,in=110] (19,0);
\end{tikzpicture}
\caption{Trzy geometrie pary zamian, narysowane jako łuki nad pozycjami.}
\label{fig:geometry}
\end{figure}

\begin{definition}
$\N_{\mathrm{dyn}}$ zawiera ruchy, których zamiany są parami rozłączne.
$\N_c=\{\sigma:\ncr(\sigma)\leq c\}$, a $\Pa_r$ zawiera ruchy,
w których każda zamiana krzyżuje się z co najwyżej $r$ innymi.
$\N_0=\Pa_0$ jest sąsiedztwem Motzkina, a $\N_\infty$ jest zbiorem
wszystkich inwolucji.
\end{definition}

„Pojedyncze krzyżowania” mogą oznaczać dwie rzeczy i badamy obie: $\N_1$
dopuszcza jedno krzyżowanie w całym ruchu, a $\Pa_1$ dopuszcza wiele par
krzyżujących się, o ile każda zamiana uczestniczy w co najwyżej jednej.
Rodziny tworzą łańcuchy
$\N_{\mathrm{dyn}}\subset\N_0\subset\N_1\subset\N_2\subset\dots\subset\N_\infty$
oraz $\N_0\subset\N_1\subset\Pa_1\subset\Pa_2\subset\dots\subset\N_\infty$,
z inkluzjami właściwymi dla dostatecznie dużych $n$; $\N_2$ i $\Pa_1$
są nieporównywalne. Rysunek~\ref{fig:geometry} pokazuje geometrie,
a Tabela~\ref{tab:taxonomy} rozmiary.

\begin{table}[t]
\centering
\caption{Rozmiar i tempo wzrostu sąsiedztw na osi nakładania się.
Rozmiary sprawdzono przez wyliczenie dla $n\leq 11$.}
\label{tab:taxonomy}
\small
\begin{tabular}{@{}lrrl@{}}
\toprule
Sąsiedztwo & $n=10$ & $n=20$ & Wzrost \\
\midrule
Dynasearch & 512 & $5{,}2{\cdot}10^5$ & $2^n$ \\
Motzkin $\N_0$ & 2\,188 & $5{,}1{\cdot}10^7$ & $3^n n^{-3/2}$ \\
$\N_1$ & 5\,038 & $2{,}6{\cdot}10^8$ & $M_n\,n/3$ \\
$\N_2$ & 7\,108 & $7{,}6{\cdot}10^8$ & $M_n\,(n/3)^2/2$ \\
$\Pa_1$ & 5\,263 & $4{,}9{\cdot}10^8$ & $3{,}4626^n$ \\
Inwolucje & 9\,496 & $2{,}4{\cdot}10^{10}$ & $(n/e)^{n/2}e^{\sqrt n}$ \\
\bottomrule
\end{tabular}
\end{table}

\begin{proposition}\label{prop:jfrac}
Niech $I_n(q)=\sum_{\sigma}q^{\ncr(\sigma)}$ po wszystkich inwolucjach
$n$ pozycji oraz $[h]_q=1+q+\dots+q^{h-1}$. Wtedy
\begin{equation}
\sum_{n\geq0} I_n(q)\,x^n=
\cfrac{1}{1-x-\cfrac{[1]_q\,x^2}{1-x-\cfrac{[2]_q\,x^2}{1-x-\cdots}}}.
\end{equation}
\end{proposition}
\begin{proof}[Szkic]
Czytamy pozycje od lewej do prawej jako ścieżkę Motzkina: końcowi
otwierającemu odpowiada krok w górę, końcowi zamykającemu krok w dół,
punktowi stałemu krok poziomy. Krok w dół z wysokości $h$ może zamknąć
każdy z $h$ otwartych łuków, a zamknięcie łuku otwartego jako $j$-ty od
końca tworzy $j$ krzyżowań. Waga kroku wynosi zatem $[h]_q$,
a twierdzenie Flajoleta~\cite{flajolet1980} daje ułamek.
\end{proof}

Dla $q=0$ otrzymujemy liczby Motzkina $M_n$, a dla $q=1$ liczby inwolucji.
Współczynnik przy $q^c$ liczy ruchy z dokładnie $c$ krzyżowaniami.

\begin{lemma}\label{lem:poisson}
Dla ustalonego $c$ liczba inwolucji $n$ pozycji z dokładnie $c$
krzyżowaniami wynosi $M_n\,(n/3)^c/c!\,(1+O(1/n))$. Stąd
$|\N_c|\sim M_n(n/3)^c/c!$, a $\N_c$ rośnie jak $3^n$ dla każdego
ustalonego $c$.
\end{lemma}
\begin{proof}[Szkic]
Na mocy Stwierdzenia~\ref{prop:jfrac} współczynnik przy $q^c$ jest sumą
po ścieżkach Motzkina z liczby sposobów rozdzielenia $c$ krzyżowań na
kroki w dół, gdzie krok z wysokości $h$ niesie co najwyżej $h-1$.
Typowa ścieżka ma $D\approx n/3$ kroków w dół, z których ograniczona
liczba zaczyna się na wysokości~1, więc liczba ta wynosi
$\mathbb{E}\binom{D+c-1}{c}=(n/3)^c/c!\,(1+O(1/n))$ na ścieżkę.
\todo{pełny dowód; sprawdzić, czy to twierdzenie jest znane dla inwolucji,
odpowiednikiem dla skojarzeń doskonałych jest $\binom{2n}{n-2}$ Riordana}
\end{proof}

Liczbowo stosunek dokładnej liczby do wzoru dla $c=1$ wynosi $0{,}82$ przy
$n=50$ i $0{,}995$ przy $n=2000$. Praktyczny wniosek jest taki, że kilka
krzyżowań niewiele dokłada: przy $n=50$ sąsiedztwo $\N_4$ obejmuje wciąż
tylko ułamek $3\cdot10^{-10}$ wszystkich inwolucji. Jednostajnie losowa
inwolucja ma około $n^2/24$ krzyżowań, więc $\N_c$ dla stałego $c$ jest
narożnikiem zbioru inwolucji, w którym prawie nie ma krzyżowań.

\begin{proposition}\label{prop:pa1}
Funkcja tworząca $S(x)$ dla $|\Pa_1|$ spełnia
$S=1+xS+x^2S^2+x^4S^4$, a $|\Pa_1|$ rośnie jak $\rho^{-n}$ z
$\rho^{-1}=3{,}4626$.
\end{proposition}
\begin{proof}[Szkic]
Pierwsza pozycja jest punktem stałym, lewym końcem łuku bez krzyżowań
albo lewym końcem łuku skrzyżowanego z dokładnie jednym innym. W ostatnim
przypadku te dwa łuki mają cztery końce, a trzy przerwy między nimi oraz
reszta po nich zawierają każda niezależną strukturę $\Pa_1$.
\end{proof}

\begin{proposition}\label{prop:optima}
Każda z powyższych rodzin zawiera wszystkie pojedyncze zamiany. Przy
wyborze addytywnym permutacja jest optimum lokalnym rodziny wtedy
i tylko wtedy, gdy jest optimum lokalnym sąsiedztwa pojedynczej zamiany.
\end{proposition}
\begin{proof}
Jeśli żadna pojedyncza zamiana nie ma ujemnej delty, to żadna suma delt nie
jest ujemna. Implikacja odwrotna zachodzi, ponieważ pojedyncze zamiany są
ruchami rodziny.
\end{proof}

Stwierdzenie~\ref{prop:optima} odróżnia oś nakładania się od osi, które
ograniczają zasięg zamiany. Nakładanie się zmienia to, jak daleko sięga
krok poza optima lokalne, ale nie zmienia zbioru optimów lokalnych.
Sąsiedztwo pojedynczej zamiany jest więc naturalną kontrolą dla całej
rodziny.

% ---------------------------------------------------------------------
\section{Dokładny wybór ruchu}\label{sec:selection}

Dla każdej zamiany-kandydata $a=(i,j)$ obliczamy
$\delta_a=C_{\max}(\pi\circ s(i,j))-C_{\max}(\pi)$ za pomocą macierzy
głów i ogonów w koszcie $O(m(j-i))$, czyli $O(mn^3)$ dla wszystkich par.
Kandydaci przechodzą filtr granic bloków z wcześniejszych
prac~\cite{windowed_own}. Przewidywana zmiana ruchu złożonego wynosi
$\sum_{a\in\sigma}\delta_a$, a wybór minimalizuje ją w obrębie rodziny.

\paragraph{Przycinanie.} Każda rodzina jest dziedziczna: podzbiór ruchu
dopuszczalnego jest dopuszczalny. Ruch optymalny używa więc tylko
kandydatów z $\delta_a<0$, a pozycje, które nie są końcem takiego
kandydata, można pominąć. Wybór działa na zbiorze $C$ kandydatów
poprawiających i na ich $E\leq 2|C|$ skompresowanych końcach.

\paragraph{Ruchy rozłączne i niekrzyżujące się.} Dla Dynasearcha wystarcza
program pozycyjny,
$B[p]=\min\bigl(B[p-1],\min_{i<p}B[i-1]+\delta_{ip}\bigr)$. Dla $\N_0$
niech $W[i][j]$ będzie najlepszą przewidywaną zmianą ruchu bez krzyżowań
wewnątrz pozycji $i..j$. Wtedy
\begin{multline}\label{eq:interval}
W[i][j]=\min\Bigl(W[i{+}1][j],\\ \min_{(i,k)\in C,\,k\leq j}
\delta_{ik}+W[i{+}1][k{-}1]+W[k{+}1][j]\Bigr),
\end{multline}
co zajmuje czas $O(E^2+E|C|)$.

\paragraph{Pojedyncze krzyżowania.} Para krzyżująca się $(i,b),(a,d)$
z $i<a<b<d$ tworzy rdzeń z dwóch łuków z trzema przerwami wewnętrznymi.
Dla $\Pa_1$ rekurencja~\eqref{eq:interval} zyskuje na pozycji $i$ jeszcze
jedną możliwość:
\begin{multline}
\delta_{ib}+\delta_{ad}+W[i{+}1][a{-}1]+W[a{+}1][b{-}1]\\
+W[b{+}1][d{-}1]+W[d{+}1][j],
\end{multline}
minimalizowaną po rdzeniach, które zaczynają się w $i$ i kończą do $j$.
Żaden łuk wewnątrz przerwy nie może krzyżować się z łukiem rdzenia,
ponieważ łuki rdzenia są już skrzyżowane jeden raz. Dla $\N_1$ program
utrzymuje dwie tablice, jedną dla ruchów bez krzyżowań i jedną dla ruchów
z dokładnie jednym rdzeniem, a rdzeń jest zamykany przerwami bez
krzyżowań. Rdzenie wylicza się z par krzyżujących się kandydatów w czasie
$O(|C|^2)$, a tablice wypełnia się w czasie $O(E^3)$. Dla $c\geq2$ spójne
składowe grafu krzyżowań mają co najwyżej $c+1$ łuków, co daje rdzenie
o ograniczonym rozmiarze i problem plecakowy po budżecie krzyżowań.
\todo{Algorytm~1: program dla $\Pa_1$ jako pseudokod; podać program dla
$c=2$ albo zostawić go dodatkowi}

\paragraph{Inwolucje.} Bez ograniczenia na krzyżowania kandydaci są
krawędziami grafu na pozycjach z wagami $-\delta_a$, a najlepszy ruch jest
skojarzeniem o maksymalnej wadze~\cite{edmonds1965}, znajdowanym w czasie
$O(n^3)$.

\paragraph{Koszt.} \todo{E0 na maszynie kampanii.} W pilotażu wybór na
stanach pobranych z trajektorii przeszukiwania zajmował co najwyżej
$0{,}6$\,ms dla $\N_0$, $4{,}5$\,ms dla $\N_1$ i $\Pa_1$ oraz $1{,}3$\,ms
dla inwolucji, wobec $14$--$221$\,ms na delty przy $n=50$--$100$.
Wyjątkiem są stany gęste. Tuż po starcie albo po ruchu, który znacznie
pogorszył uszeregowanie, poprawia się tysiące kandydatów, a wąskim
gardłem staje się wyliczanie rdzeni. Dlatego ograniczamy rdzenie do $400$
najbardziej poprawiających kandydatów, a skojarzenie do $3000$ najbardziej
poprawiających, i zapisujemy, jak często ograniczenie jest aktywne.

\paragraph{Akceptacja.} Wybrany ruch jest stosowany bez zmian; gdy żaden
kandydat nie poprawia, stosowana jest najlepsza pojedyncza zamiana. Jest to
reguła wcześniejszych implementacji Dynasearcha i Motzkina i używają jej
wszystkie wyniki badania głównego. Sekcja~\ref{sec:safe} omawia wariant,
który przed zastosowaniem ruchu sprawdza rzeczywistą długość
uszeregowania.

% ---------------------------------------------------------------------
\section{Addytywność a geometria pary}\label{sec:additivity}

Wybór optymalizuje sumę delt, ale zmiana długości uszeregowania po dwóch
zamianach nie jest sumą ich zmian. W ujęciu ścieżek kratowych długość
uszeregowania jest długością najdłuższej ścieżki przez siatkę operacji,
$C'_{\max}=\max_P\bigl[\ell(P)+\sum_a\Delta_a(P)\bigr]$. Dwie zamiany
rozłączne oddziałują ze sobą tylko przez wybór ścieżki krytycznej. Para
zagnieżdżona lub krzyżująca się zmienia także zadania leżące wewnątrz
przedziału drugiej zamiany, więc każda delta jest liczona w kontekście,
który druga zamiana niszczy. Składnik drugiego rzędu zapisujemy jako
$Q_{ab}=\Delta_{ab}-\delta_a-\delta_b$, gdzie $\Delta_{ab}$ jest
rzeczywistą zmianą pary. \todo{minimalne kontrprzykłady dla każdej
geometrii i warunek wystarczający na $Q_{ab}=0$}

\begin{table}[t]
\centering
\caption{Zrealizowany udział przewidywanego wspólnego zysku dwóch
poprawiających zamian według geometrii, średnia po parach na stanach
z trajektorii.}
\label{tab:geometry}
\small
\begin{tabular}{@{}lccc@{}}
\toprule
Klasa & rozłączne & zagnieżdżone & krzyżujące się \\
\midrule
$20\times10$ & 0,80 & 0,65 & 0,58 \\
$20\times20$ & 0,72 & 0,60 & 0,46 \\
$50\times10$ & 0,90 & 0,85 & 0,83 \\
$50\times20$ & 0,82 & 0,76 & 0,71 \\
$100\times10$ & 0,94 & 0,90 & 0,90 \\
$100\times20$ & 0,87 & 0,81 & 0,79 \\
$200\times20$ & 0,92 & 0,90 & 0,89 \\
\bottomrule
\end{tabular}
\end{table}

Pobieramy około 2000 stanów z trajektorii ILS z kampanii, około 190 na
klasę, i mierzymy każdą parę poprawiających zamian-kandydatów w każdym
stanie. Tabela~\ref{tab:geometry} pokazuje wynik. W każdej klasie pary
rozłączne zachowują większą część przewidywanego zysku niż pary nakładające
się (test Wilcoxona w każdej klasie, $p<10^{-7}$). Różnica jest największa
dla $n=20$: dla $20\times20$ pary rozłączne zachowują $0{,}72$, zagnieżdżone
$0{,}60$, a krzyżujące się $0{,}46$. Pary krzyżujące się nigdy nie są lepsze
od zagnieżdżonych, a dla $n\geq50$ obie są prawie na tym samym poziomie.
Nakładanie się psuje addytywność, a krzyżowanie niewiele dodaje poza tym.

Dla całych ruchów efekt się kumuluje. Ruchy Dynasearcha realizują
$0{,}46$--$0{,}88$ swojego przewidywanego zysku, ruchy z $\N_0$, $\N_1$
i $\Pa_1$ tylko $0{,}03$--$0{,}21$, a najlepsza inwolucja pogarsza średnio
długość uszeregowania w każdej klasie. Wybrany ruch złożony jest gorszy od
najlepszej pojedynczej zamiany w 11--44\% stanów dla Dynasearcha,
w 48--71\% dla $\N_0$, $\N_1$ i $\Pa_1$ oraz w 74--91\% dla inwolucji.
Optimum inwolucyjne wymaga także wielu krzyżowań. Najmniejsze $c$,
którego optimum je osiąga, wynosi co najmniej dwa w 65--100\% stanów,
a na 60 najlepszych kandydatach budżet dwóch krzyżowań osiąga je w co
najwyżej 40\% stanów przy $n=20$, w co najwyżej 11\% przy $n=100$ i nigdy
przy $n=200$.

Większa rodzina wybiera też więcej zamian na ruch, z medianą 3 dla
Dynasearcha, 15 dla $\N_0$ i 21 dla inwolucji, więc geometria i liczność
są pomieszane. Przy równej liczności porządek się utrzymuje. Ruchy
z dwiema zamianami realizują medianę $0{,}79$ swojej prognozy
w Dynasearchu, $0{,}38$ w $\N_0$, $0{,}30$ w $\Pa_1$ i $0{,}29$ dla
inwolucji, a ruchy z trzema lub czterema zamianami $0{,}70$, $0{,}20$,
$-0{,}07$ i $-0{,}19$. Oba czynniki mają znaczenie. Każda rodzina wybiera
inne ruchy, więc część tej różnicy może pochodzić z samego wyboru;
Sekcja~\ref{sec:exp} kontroluje zatem liczność bezpośrednio.
\todo{rysunek: przewidywany i zrealizowany zysk wzdłuż osi}

% ---------------------------------------------------------------------
\section{Eksperymenty}\label{sec:exp}

Protokół jest taki jak we wcześniejszych pracach~\cite{windowed_own},
a rodzina sąsiedztw jest jedynym czynnikiem. Instancje to 120 instancji
Taillarda~\cite{taillard1993} w 12 klasach z $n\in\{20,50,100,200,500\}$
i $m\in\{5,10,20\}$. Metaheurystyką jest niezmienione ILS z listą tabu
o długości 10 i pulą elitarną 10 najlepszych rozwiązań zaburzanych
kopnięciami typu podwójny most (double-bridge), startujące z permutacji
tożsamościowej. Każda konfiguracja jest uruchamiana jednokrotnie przez
$10$\,s na ziarno, z 5 ziarnami, a najlepsze rozwiązania w chwilach $0{,}1$,
$0{,}5$, $1$, $2$ i $5$\,s są odczytywane z dziennika trajektorii, co jest
poprawne, ponieważ limit czasu działa tylko jako reguła zatrzymania.
Warianty to Dynasearch, $\N_0$, $\N_1$, $\Pa_1$ i inwolucje, z najlepszą
pojedynczą zamianą jako kontrolą. Wszystkie warianty używają jednej
implementacji kandydatów, delt i filtra. Uruchomienia są sekwencyjne na
Apple M4 Pro. Surowe przebiegi i skrypty odtwarzające każdą tabelę
i rysunek zostaną udostępnione po zakończeniu recenzji.

Jakością jest względne odchylenie procentowe (RPD) od dolnego ograniczenia
Taillarda po $10$\,s. Wnioski wyciągamy tylko w klasach, w których
$10$-sekundowy przebieg wykonuje co najmniej 30 ruchów, czego oczekujemy
dla $n\leq100$ (dziewięć klas); przy $n=500$ przebieg wykonuje jeden ruch.
Dwie rodziny testów ustalono przed kampanią, każdą z instancjami jako
blokami i z poprawką Holma po klasach: inwolucje wobec $\N_0$ oraz
Dynasearch wobec $\N_0$, testy dwustronne. Test trendu Page'a~\cite{page1963}
stosujemy wzdłuż łańcucha
$\N_{\mathrm{dyn}}\subset\N_0\subset\N_1\subset\Pa_1\subset\N_\infty$.
Wszystkie pozostałe porównania są eksploracyjne i tak je podajemy. Rodziny
z krzyżowaniami wykonują mniej iteracji, więc pierwszą rodzinę testów
analizujemy także przy równej liczbie iteracji. Kontrole z drugiego etapu
używają tego samego protokołu na dziewięciu klasach z $n\leq100$,
a kontrola QUBO na sześciu klasach z $n\leq50$.

\begin{table*}[t]
\centering
\caption{Średnie RPD [\%] do dolnego ograniczenia po $10$\,s, ILS, według
klas (10 instancji, 5 ziaren).}
\label{tab:main}
\small
\begin{tabular}{@{}l*{12}{c}@{}}
\toprule
Wariant & \multicolumn{3}{c}{$n=20$} & \multicolumn{3}{c}{$n=50$} &
\multicolumn{3}{c}{$n=100$} & \multicolumn{2}{c}{$n=200$} & $n=500$ \\
$m$ & 5 & 10 & 20 & 5 & 10 & 20 & 5 & 10 & 20 & 10 & 20 & 20 \\
\midrule
Pojedyncza zamiana & 2,61 & 9,98 & 21,36 & 1,10 & 5,01 & 13,92 & 1,46 & 2,74 & 10,24 & 5,61 & 17,37 & 15,76 \\
Dynasearch & 2,52 & 9,80 & 21,29 & 1,01 & 4,57 & 13,33 & 1,36 & 2,46 & 10,34 & 4,40 & 15,62 & 14,49 \\
$\N_0$ (Motzkin) & 2,62 & 10,19 & 21,44 & 2,00 & 12,60 & 19,56 & 4,13 & 12,19 & 20,39 & 12,17 & 19,70 & 14,65 \\
$\N_1$ & 2,66 & 10,30 & 21,56 & 2,43 & 12,09 & 21,10 & 7,05 & 12,84 & 20,15 & 12,98 & 19,67 & 14,65 \\
$\Pa_1$ & 2,67 & 10,48 & 21,62 & 2,78 & 13,01 & 21,15 & 6,61 & 12,00 & 19,47 & 12,88 & 19,12 & 14,64 \\
Inwolucje & 3,31 & 12,48 & 23,16 & 6,05 & 20,91 & 26,27 & 9,23 & 18,26 & 25,70 & 15,90 & 22,53 & 15,99 \\
\bottomrule
\end{tabular}
\end{table*}

\paragraph{Wynik główny.} Tabela~\ref{tab:main} pokazuje kampanię.
Jakość maleje wzdłuż osi, a test trendu Page'a wzdłuż
$\N_{\mathrm{dyn}}\subset\N_0\subset\N_1\subset\Pa_1\subset\N_\infty$
jest istotny w każdej klasie (łącznie dla $n\leq100$, $p\approx10^{-53}$).
W pierwszej zarejestrowanej rodzinie testów inwolucje są gorsze od
$\N_0$ we wszystkich 12 klasach, we wszystkich dziewięciu rozstrzygających
po poprawce Holma, a także przy równej liczbie iteracji, więc strata nie
wynika z wolniejszych ruchów. W drugiej rodzinie Dynasearch jest lepszy od
$\N_0$ w siedmiu z dziewięciu klas rozstrzygających, a dla $20\times5$
i $20\times20$ oba się nie różnią. Spadek jest nierówny. $\N_1$ i $\Pa_1$
pozostają blisko $\N_0$, a dwa duże kroki prowadzą od Dynasearcha do
$\N_0$, gdzie wchodzi zagnieżdżanie, oraz od $\Pa_1$ do inwolucji, gdzie
krzyżowania przestają być ograniczone. Dla $n=50$ i $n=100$ rodziny
nakładające się przegrywają z Dynasearchem o 1--10 punktów, a inwolucje
o nawet 16. Rodziny z krzyżowaniami wykonują też mniej ruchów, od około
połowy ruchów Dynasearcha dla $20\times5$ do jednej czwartej dla
$100\times5$.

Ruchy złożone dodają niewiele do najlepszej pojedynczej zamiany.
Dynasearch wyprzedza ją o $0{,}1$--$0{,}6$ punktu dla $n\leq100$ i o
$1{,}2$--$1{,}8$ punktu dla $n=200$, a po poprawce Holma istotna jest
tylko ta druga różnica. Dla $n=200$ przebieg wykonuje 6--15 ruchów,
a dla $n=500$ jeden ruch, więc te klasy nic nie mówią o $c$.

Te wzorce zgadzają się z pomiarem z Sekcji~\ref{sec:additivity}. ILS
z listą tabu stosuje wybrany ruch nawet wtedy, gdy pogarsza on
uszeregowanie. Dla zamian rozłącznych taki ruch jest nieznacznie gorszy;
dla zamian nakładających się często jest znacznie gorszy. Ten sam
mechanizm tłumaczy słabe wyniki Motzkina dla $n\geq50$ we wcześniejszym
badaniu skalowania~\cite{windowed_own}, które stosowało tę samą regułę.
\todo{sformułowanie uzgodnione z promotorem}

\paragraph{Kontrole.} Trzy kontrole sprawdzają to wyjaśnienie. Po
pierwsze, tniemy każdy ruch do jego $K$ najbardziej poprawiających zamian.
Dla $K=2$ wszystkie rodziny działają podobnie: spośród 27 porównań
parami między Dynasearchem, $\N_0$ i inwolucjami tylko jedno jest istotne
po poprawce Holma i tam wygrywają inwolucje. Ruchy ucięte dla $\N_0$
i dla inwolucji są lepsze od pełnych o nawet 11 i 16 punktów. Dla $K=4$
inwolucje są już gorsze od $\N_0$ w czterech klasach z $n\leq50$.
Krzyżowania szkodzą więc przez oddziaływanie wielu jednoczesnych zamian,
a nie przez pojedynczą parę krzyżującą się. Po drugie, akceptujemy
najlepszy spośród pierwszego ruchu, do siedmiu jego alternatyw
otrzymanych przez pominięcie jednej z jego zamian, oraz najlepszej
pojedynczej zamiany, wszystkie ocenione dokładnie. Ta reguła top-8 nie
jest lepsza od pojedynczego sprawdzenia z Sekcji~\ref{sec:safe} w żadnej
klasie. Po trzecie, ciaśniejsze ograniczenie kandydatów (1000 par, 100
w rdzeniu) czyni każdy ruch tańszym i poprawia $\N_0$ oraz inwolucje dla
$n\geq100$ o nawet $2{,}5$ punktu, ale pozostają one daleko za
Dynasearchem.

\paragraph{Kontekst.} Nasze ILS jest punktem odniesienia, a nie
konkurencyjnym solwerem problemu przepływowego. Uruchomione z NEH~\cite{nawaz1983}
zamiast z permutacji tożsamościowej zyskuje do $1{,}9$ punktu dla
$n=100$, a Dynasearch jest wtedy w granicach $0{,}13$ punktu od pojedynczej
zamiany dla każdego $n\leq100$. Iterowany algorytm zachłanny~\cite{ruiz2007},
uruchomiony w tej samej bazie kodu z tym samym budżetem, jest lepszy od
obu w siedmiu z dziewięciu klas po poprawce Holma, o nawet $2{,}6$ punktu.
Nasze porównanie dotyczy więc ruchów złożonych z zamian wewnątrz jednej
metaheurystyki, a nie stanu wiedzy dla problemu przepływowego.

% ---------------------------------------------------------------------
\section{Wariant z bezpiecznikiem}\label{sec:safe}

Pilotaż nasuwa oczywiste pytanie: jak dużą część straty dają ruchy, które
przeszukiwanie by odrzuciło, gdyby spojrzało na ich rzeczywistą długość
uszeregowania? Aby odpowiedzieć, uruchamiamy osobny wariant. Długość
uszeregowania wybranego ruchu złożonego jest oceniana jednokrotnie
w koszcie $O(mn)$, a stosowany jest lepszy z tego ruchu i najlepszej
pojedynczej zamiany, więc krok nigdy nie jest gorszy od najlepszej
pojedynczej zamiany. Zmienia to regułę wyboru, a nie sąsiedztwo, więc
wariant pozostaje poza głównymi tabelami i poza zarejestrowanymi testami
i nie jest porównywalny z wcześniejszymi wynikami Dynasearcha i Motzkina.

\begin{table*}[t]
\centering
\caption{Wariant z bezpiecznikiem: średnie RPD [\%] do dolnego ograniczenia
po $10$\,s, ten sam protokół co w Tabeli~\ref{tab:main}.}
\label{tab:safe}
\small
\begin{tabular}{@{}l*{12}{c}@{}}
\toprule
Wariant & \multicolumn{3}{c}{$n=20$} & \multicolumn{3}{c}{$n=50$} &
\multicolumn{3}{c}{$n=100$} & \multicolumn{2}{c}{$n=200$} & $n=500$ \\
$m$ & 5 & 10 & 20 & 5 & 10 & 20 & 5 & 10 & 20 & 10 & 20 & 20 \\
\midrule
Dynasearch & 2,53 & 9,87 & 21,34 & 1,06 & 4,84 & 13,56 & 1,42 & 2,61 & 9,75 & 4,43 & 14,83 & 14,49 \\
$\N_0$ (Motzkin) & 2,51 & 9,92 & 21,21 & 1,00 & 4,69 & 13,54 & 1,31 & 2,58 & 9,56 & 4,61 & 14,89 & 14,55 \\
$\N_1$ & 2,49 & 9,90 & 21,22 & 1,04 & 4,82 & 13,42 & 1,33 & 2,46 & 9,45 & 4,80 & 14,92 & 14,48 \\
$\Pa_1$ & 2,49 & 9,92 & 21,25 & 1,04 & 4,92 & 13,53 & 1,37 & 2,45 & 9,59 & 4,77 & 14,76 & 14,53 \\
Inwolucje & 2,50 & 9,89 & 21,25 & 1,01 & 4,95 & 13,38 & 1,39 & 2,45 & 9,65 & 5,92 & 16,38 & 15,58 \\
\bottomrule
\end{tabular}
\end{table*}

Tabela~\ref{tab:safe} pokazuje, że sprawdzenie usuwa stratę rodzin
nakładających się. Dla $50\times10$ obniża ono $\N_0$ z $12{,}60$ do
$4{,}69$, a inwolucje z $20{,}91$ do $4{,}95$, a dla $n\leq100$ pięć
rodzin z bezpiecznikiem leży w granicach pół punktu od siebie w każdej
klasie. Dla $\N_0$, $\N_1$, $\Pa_1$ i inwolucji sprawdzenie jest istotnie
lepsze od zwykłej reguły w prawie każdej klasie z $n\geq50$. $\N_0$
z bezpiecznikiem nie różni się istotnie od zwykłego Dynasearcha w żadnej
klasie: wyprzedza go o około $0{,}75$ punktu dla $100\times20$
i $200\times20$ i ustępuje o co najwyżej $0{,}2$ w pozostałych.
Dla $n\geq200$ inwolucje z bezpiecznikiem zostają w tyle, ponieważ
działa tam ograniczenie liczby kandydatów. Tanie sprawdzenie rzeczywistej
długości uszeregowania znaczy więcej niż wybór spośród rodzin
nakładających się.

% ---------------------------------------------------------------------
\section{Struktura QUBO}\label{sec:quantum}

Jako QUBO nad zamianami-kandydatami rodzina jest zbiorem kar na parach
kandydatów. QUBO Motzkina karze pary, które mają wspólny koniec, oraz
pary krzyżujące się. QUBO inwolucji karze tylko wspólne końce. Największe
sąsiedztwo na osi ma więc najrzadsze QUBO. Trzeci model dodaje zmierzone
składniki $Q_{ab}$ i wyraża to, czego nie wyrazi żaden program po
przedziałach; jest najgęstszy.

Zanim zaangażujemy jakikolwiek procesor, testujemy trzeci model
klasycznie. Wewnątrz ILS QUBO z poprawką na pary, zbudowane nad 16
najbardziej poprawiającymi zamianami i rozwiązane dokładnie, zastępuje
wybór addytywny. Na sześciu klasach z $n\leq50$ jest lepsze od zwykłych
inwolucji o $0{,}7$--$16$ punktów i lepsze od $\N_0$ w pięciu klasach, nie
różni się od Dynasearcha i jest lepsze od najlepszej pojedynczej zamiany
w dwóch klasach. Na stanach z analizy anatomii jego optimum realizuje
medianę $0{,}85$--$1{,}00$ swojego przewidywanego zysku i jest gorsze od
najlepszej pojedynczej zamiany w co najwyżej 10\% stanów, wobec 11--71\%
dla optimum inwolucyjnego nad tymi samymi 16 kandydatami. Składniki par
naprawiają to, co psuje addytywność, i są to składniki, których program po
przedziałach nie wyrazi.

\begin{table}[t]
\centering
\caption{PILOTAŻ: sprzężenia i fizyczne kubity na grafie Pegasus dla QUBO
Motzkina (M), inwolucji (I) i z poprawką na pary (PQ) nad tymi samymi $K$
kandydatami.}
\label{tab:qubo}
\small
\begin{tabular}{@{}lrcc@{}}
\toprule
Stan & $K$ & Sprzężenia M / I / PQ & Kubity M / I / PQ \\
\midrule
$20\times10$ & 18 & 59 / 41 / 140 & 26 / 22 / 45 \\
$20\times20$ & 11 & 27 / 13 / 55 & 14 / 10 / 20 \\
$50\times10$ & 105 & 2117 / 634 / -- & 1001 / 425 / -- \\
\bottomrule
\end{tabular}
\end{table}

Tabela~\ref{tab:qubo} pokazuje strukturę na trzech stanach z trajektorii.
Planujemy użyć procesora jako wyroczni na zapisanych stanach: dla każdego
stanu trzy QUBO nad tymi samymi kandydatami są próbkowane i porównywane
z optimum dokładnym. Rozróżniamy dwa punkty końcowe. Wierność sprzętu to
odstęp energii od optimum próbkowanego modelu i udział próbek
dopuszczalnych. Jakość ruchu to rzeczywista zmiana długości uszeregowania.
QUBO inwolucji może wygrać w pierwszym i przegrać w drugim, ponieważ jego
optimum zawiera więcej zamian nakładających się. \todo{E4: wyniki
annealera i komputera bramkowego na zapisanych stanach}

% ---------------------------------------------------------------------
\section{Wnioski}
Przy wyborze addytywnym zamiany ruchu złożonego nie powinny się
krzyżować, a w tym protokole nie powinny nawet się zagnieżdżać. Jakość
maleje wzdłuż osi od Dynasearcha do inwolucji w każdej klasie Taillarda,
a większe sąsiedztwa więcej przewidują i mniej dostarczają. Przyczyną
jest oddziaływanie wielu nakładających się zamian. Ucięte do dwóch zamian
na ruch wszystkie rodziny działają podobnie, a przy zmierzonym modelu
oddziaływań par, rozwiązanym dokładnie jako QUBO, ruchy z krzyżowaniami
dorównują Dynasearchowi. Jednokrotne sprawdzenie rzeczywistej długości
uszeregowania daje ten sam efekt przy niższym koszcie. Ruchy złożone
z zamian dowolnego kształtu dodają w tym ILS niewiele do najlepszej
pojedynczej zamiany, a iterowany algorytm zachłanny jest znacznie przed
wszystkimi z nich. Otwarte pozostaje pytanie, czy model z poprawką na
pary nadal się opłaca, gdy zbiór kandydatów przerasta rozwiązanie
dokładne; w tym miejscu pojawiają się annealery i inne samplery QUBO.

% ---------------------------------------------------------------------
% All entries to be verified in Crossref (authors, volume, pages, DOI).
\begin{thebibliography}{99}\small
\bibitem{ahuja2002} R.K.~Ahuja, \"O.~Ergun, J.B.~Orlin, A.P.~Punnen.
A survey of very large-scale neighborhood search techniques.
\emph{Discrete Appl.\ Math.} 123:75--102, 2002.
\bibitem{akutsu2000} T.~Akutsu. Dynamic programming algorithms for RNA
secondary structure prediction with pseudoknots. \emph{Discrete Appl.\
Math.} 104:45--62, 2000. \todo{verify}
\bibitem{bozejko2007} W.~Bo\.zejko, M.~Wodecki. On the theoretical
properties of swap multimoves. \emph{Oper.\ Res.\ Lett.} 35(2):227--231, 2007.
\bibitem{bozejko2026sr} W.~Bo\.zejko, M.~Uchro\'nski, M.~Wodecki. The new
exponential Fibonacci neighborhood as QUBO generated on a D-Wave quantum
computing environment for the permutational scheduling problem.
\emph{Sci.\ Rep.}, 2026.
\bibitem{chen2007} W.Y.C.~Chen, E.Y.P.~Deng, R.R.X.~Du, R.P.~Stanley,
C.H.~Yan. Crossings and nestings of matchings and partitions.
\emph{Trans.\ Amer.\ Math.\ Soc.} 359:1555--1575, 2007. \todo{verify}
\bibitem{congram2002} R.K.~Congram, C.N.~Potts, S.L.~van de Velde. An
iterated dynasearch algorithm for the single-machine total weighted
tardiness scheduling problem. \emph{INFORMS J.\ Comput.} 14(1):52--67, 2002.
\bibitem{deineko2000} V.G.~Deineko, G.J.~Woeginger. A study of
exponential neighborhoods for the travelling salesman problem and for
the quadratic assignment problem. \emph{Math.\ Program.} 87:519--542,
2000. \todo{verify}
\bibitem{edmonds1965} J.~Edmonds. Paths, trees, and flowers.
\emph{Canad.\ J.\ Math.} 17:449--467, 1965.
\bibitem{ergun2006} \"O.~Ergun, J.B.~Orlin. Fast neighborhood search for
the single machine total weighted tardiness problem. \emph{Oper.\ Res.\
Lett.} 34(1):41--45, 2006. \todo{verify}
\bibitem{fibonacci_own} \todo{own work, third person: Fibonacci neighborhood}
\bibitem{flajolet1980} P.~Flajolet. Combinatorial aspects of continued
fractions. \emph{Discrete Math.} 32(2):125--161, 1980.
\bibitem{garey1976} M.R.~Garey, D.S.~Johnson, R.~Sethi. The complexity of
flowshop and jobshop scheduling. \emph{Math.\ Oper.\ Res.} 1(2):117--129, 1976.
\bibitem{grosso2004} A.~Grosso, F.~Della Croce, R.~Tadei. An enhanced
dynasearch neighborhood for the single-machine total weighted tardiness
scheduling problem. \emph{Oper.\ Res.\ Lett.} 32(1):68--72, 2004. \todo{verify}
\bibitem{gutin1999} G.~Gutin. Exponential neighbourhood local search for
the traveling salesman problem. \emph{Comput.\ Oper.\ Res.}
26:313--320, 1999. \todo{verify}
\bibitem{kasraoui2006} A.~Kasraoui, J.~Zeng. Distribution of crossings,
nestings and alignments of two edges in matchings and partitions.
\emph{Electron.\ J.\ Combin.} 13, R33, 2006. \todo{verify}
\bibitem{lyngso2000} R.B.~Lyngs\o, C.N.S.~Pedersen. RNA pseudoknot
prediction in energy-based models. \emph{J.\ Comput.\ Biol.}
7(3--4):409--427, 2000. \todo{verify}
\bibitem{motzkin_own} \todo{own work, third person: Motzkin neighborhood}
\bibitem{nussinov1978} R.~Nussinov, G.~Pieczenik, J.R.~Griggs,
D.J.~Kleitman. Algorithms for loop matchings. \emph{SIAM J.\ Appl.\
Math.} 35(1):68--82, 1978.
\bibitem{nawaz1983} M.~Nawaz, E.E.~Enscore, I.~Ham. A heuristic algorithm
for the $m$-machine, $n$-job flow-shop sequencing problem. \emph{Omega}
11(1):91--95, 1983.
\bibitem{page1963} E.B.~Page. Ordered hypotheses for multiple
treatments: a significance test for linear ranks. \emph{J.\ Amer.\
Statist.\ Assoc.} 58:216--230, 1963.
\bibitem{riordan1975} J.~Riordan. The distribution of crossings of
chords joining pairs of $2n$ points on a circle. \emph{Math.\ Comp.}
29:215--222, 1975.
\bibitem{rivas1999} E.~Rivas, S.R.~Eddy. A dynamic programming algorithm
for RNA structure prediction including pseudoknots. \emph{J.\ Mol.\
Biol.} 285(5):2053--2068, 1999.
\bibitem{ruiz2007} R.~Ruiz, T.~St\"utzle. A simple and effective iterated
greedy algorithm for the permutation flowshop scheduling problem.
\emph{Eur.\ J.\ Oper.\ Res.} 177(3):2033--2049, 2007.
\bibitem{sarvanov1981} V.I.~Sarvanov, N.N.~Doroshko. Approximate
solution of the travelling salesman problem by a local algorithm that
searches neighbourhoods of factorial cardinality in cubic time.
1981. \todo{verify source}
\bibitem{schuirmann1987} D.J.~Schuirmann. A comparison of the two
one-sided tests procedure and the power approach for assessing the
equivalence of average bioavailability. \emph{J.\ Pharmacokinet.\
Biopharm.} 15:657--680, 1987.
\bibitem{tabaska1998} J.E.~Tabaska, R.B.~Cary, H.N.~Gabow, G.D.~Stormo.
An RNA folding method capable of identifying pseudoknots and base
triples. \emph{Bioinformatics} 14(8):691--699, 1998.
\bibitem{taillard1993} E.~Taillard. Benchmarks for basic scheduling
problems. \emph{Eur.\ J.\ Oper.\ Res.} 64(2):278--285, 1993.
\bibitem{touchard1952} J.~Touchard. Sur un probl\`eme de configurations
et sur les fractions continues. \emph{Canad.\ J.\ Math.} 4:2--25, 1952.
\bibitem{windowed_own} \todo{own work, third person: windowed QUBO / classical--quantum neighborhoods}
\end{thebibliography}
\end{document}
